\section*{Capítulo \ref{ch:b2:diatomic-mos} --- Orbitais moleculares das moléculas diatômicas}

\begin{solution}{exo:b2:diatomic-mos:1}
$E_+ = (-13.6 - 9.0)/1.60 = \qty{-14.1}{eV}$; $E_- = (-13.6 + 9.0)/0.40 = \qty{-11.5}{eV}$.
Dois elétrons em $\psi_+$: $-28.25$ contra $2 \times (-13.6) = \qty{-27.2}{eV}$, ligada por
\qty{1.05}{eV} nesse modelo grosseiro.
\end{solution}

\begin{solution}{exo:b2:diatomic-mos:2}
\ce{He2}: $(2 - 2)/2 = 0$; \ce{He2+}: $(2 - 1)/2 = 0.5$; \ce{Li2}: 1; \ce{B2}: seis
elétrons de valência, $\sigma_{2s}^2\sigma_{2s}^{*2}\pi_{2p}^2$, \omterm{def:b2:diatomic-mos:bond-order}{ordem de ligação} 1.
\end{solution}

\begin{solution}{exo:b2:diatomic-mos:3}
\ce{B2} (dois elétrons desemparelhados nos dois orbitais $\pi_{2p}$, com mistura s--p) e
\ce{O2} (dois em $\pi^*$). \ce{C2}, \ce{N2}, \ce{F2} são \omterm{def:b2:diatomic-mos:magnetism}{diamagnéticas}.
\end{solution}

\begin{solution}{exo:b2:diatomic-mos:4}
Não nula: (1s, $2\mathrm p_z$), ($2\mathrm p_x$, $2\mathrm p_x$), (2s, $2\mathrm p_z$).
Nula por simetria: (1s, $2\mathrm p_x$) e ($2\mathrm p_x$, $2\mathrm p_y$).
\end{solution}

\begin{solution}{exo:b2:diatomic-mos:5}
A ionização retira um elétron ligante: \omterm{def:b2:diatomic-mos:bond-order}{ordem de ligação} 2.5 em vez de 3, logo uma
ligação mais longa e mais fraca, como mostra $D_0(\ce{N2+}) = \qty{8.71}{eV} < D_0(\ce{N2}) = \qty{9.76}{eV}$.
\end{solution}

\begin{solution}{exo:b2:diatomic-mos:6}
Os orbitais do oxigênio ficam mais baixos (o oxigênio é mais eletronegativo):
pelo \cref{thm:b2:diatomic-mos:heteronuclear}, os \omterm{def:b2:diatomic-mos:bonding}{orbitais ligantes} têm maior
peso no átomo mais baixo, o oxigênio, e os orbitais mais altos, no carbono. O
orbital ocupado mais alto, um orbital $\sigma$ de caráter sobretudo carbônico
voltado para o lado oposto ao oxigênio, é o par não ligante do carbono.
\end{solution}

\begin{solution}{exo:b2:diatomic-mos:7}
Só o $2\mathrm p_z$ do flúor (ao longo do eixo) tem a simetria do 1s do
hidrogênio: eles formam um orbital $\sigma$ com maior peso em F e um $\sigma^*$
com maior peso em H. Os $2\mathrm p_x$, $2\mathrm p_y$ do flúor (sobreposição
nula) e seu 2s baixo (distante demais em energia) permanecem não ligantes:
três pares não ligantes. Os oito elétrons de valência preenchem o 2s, o
$\sigma$ e os dois p não ligantes.
\end{solution}

\begin{solution}{exo:b2:diatomic-mos:8}
$\bar\alpha = \qty{-12}{eV}$, $\Delta = \qty{4}{eV}$, $\sqrt{4 + 9} = \qty{3.61}{eV}$:
$E = \qty{-15.6}{eV}$ e \qty{-8.4}{eV}. Para o nível mais baixo, $(\alpha_B - E)c_B = -\beta
c_A$: $1.61c_B = 3.0c_A$, $c_A/c_B = 0.535$, $c_B^2 = 1/(1 + 0.286) = 0.78$.
\end{solution}

\begin{solution}{exo:b2:diatomic-mos:9}
Oito elétrons de valência: $\sigma_{2s}^2\sigma_{2s}^{*2}\pi_{2p}^4$, estando vazio o orbital
$\sigma_{2p}$ (empurrado para cima de $\pi_{2p}$ pela mistura). \omterm{def:b2:diatomic-mos:bond-order}{Ordem de ligação}
$(6 - 2)/2 = 2$, devida aos dois orbitais $\pi$ preenchidos.
\end{solution}

\begin{solution}{exo:b2:diatomic-mos:10}
NO, onze elétrons de valência, um em $\pi^*$: \omterm{def:b2:diatomic-mos:bond-order}{ordem de ligação} 2.5. \ce{NO+}, dez: \omterm{def:b2:diatomic-mos:bond-order}{ordem
de ligação} 3. $D_0(\ce{NO+}) = 6.51 + 13.62 - 9.26 = \qty{10.87}{eV}$: muito mais forte
que NO, pois o elétron retirado era antiligante.
\end{solution}

\begin{solution}{exo:b2:diatomic-mos:11}
$E_+ + E_- = \dfrac{(\alpha + \beta)(1 - S) + (\alpha - \beta)(1 + S)}{1 - S^2} = \dfrac{2(\alpha -
\beta S)}{1 - S^2}$. Ela supera $2\alpha$ quando $\alpha - \beta S > \alpha - \alpha S^2$, isto é,
$\beta < \alpha S$ (dividindo por $-S < 0$). Com quatro elétrons, dois em cada orbital,
a energia total fica acima da dos orbitais separados: dois orbitais preenchidos
se repelem.
\end{solution}

\begin{solution}{exo:b2:diatomic-mos:12}
\omterm{def:b2:diatomic-mos:bond-order}{Ordens de ligação} 2.5, 2, 1.5, 1: nessa ordem, as ligações se alongam e se
enfraquecem. O peróxido de hidrogênio contém a ligação simples O--O do
peróxido (unidade \ce{O2^2-}); \ce{KO2} contém o íon superóxido \ce{O2-},
\omterm{def:b2:diatomic-mos:magnetism}{paramagnético}.
\end{solution}

\begin{solution}{pb:b2:diatomic-mos:1}
\textbf{1.} Doze.
\textbf{2.} $\sigma_{2s}^2\,\sigma_{2s}^{*2}\,\sigma_{2p}^2\,\pi_{2p}^4\,\pi_{2p}^{*2}$, com os dois elétrons $\pi^*$
em orbitais diferentes.
\textbf{3.} $(8 - 4)/2 = 2$.
\textbf{4.} Dois: \ce{O2} é \omterm{def:b2:diatomic-mos:magnetism}{paramagnético}.
\textbf{5.} Ela emparelha todos os elétrons, enquanto dois deles ficam sozinhos
em dois orbitais $\pi^*$ degenerados.
\textbf{6.} $D_0 = 2 \times 246.79 = \qty{493.6}{kJ/mol}$, \qty{5.12}{eV}.
\textbf{7.} \ce{O2+}: $\pi^{*1}$, \omterm{def:b2:diatomic-mos:bond-order}{ordem de ligação} 2.5.
\textbf{8.} \ce{O2-}: $\pi^{*3}$, \omterm{def:b2:diatomic-mos:bond-order}{ordem de ligação} 1.5.
\textbf{9.} \ce{O2^2-}: $\pi^{*4}$, \omterm{def:b2:diatomic-mos:bond-order}{ordem de ligação} 1.
\textbf{10.} \ce{O2+} um, \ce{O2-} um, \ce{O2^2-} nenhum.
\textbf{11.} \ce{O2+} $<$ \ce{O2} $<$ \ce{O2-} $<$ \ce{O2^2-}.
\textbf{12.} O íon peróxido.
\textbf{13.} \ce{O2+}, \ce{O2} e \ce{O2-}.
\textbf{14.} De um orbital $\pi^*$.
\textbf{15.} O orbital $\pi^*$, antiligante, fica acima do nível 2p do átomo: seu
elétron está menos ligado.
\textbf{16.} Dois caminhos de \ce{O2} a $\ce{O} + \ce{O+} + \mathrm e^-$: dissociar e depois
ionizar um átomo, $D_0(\ce{O2}) + I(\ce{O})$; ou ionizar a molécula e depois dissociar
o íon, $I(\ce{O2}) + D_0(\ce{O2+})$. Iguais: $D_0(\ce{O2+}) = 5.12 + 13.62 - 12.07 = \qty{6.66}{eV}$.
\textbf{17.} Maior que $D_0(\ce{O2}) = \qty{5.12}{eV}$: retirar um elétron antiligante
fortalece a ligação (ordem 2.5).
\textbf{18.} $D_0(\ce{N2}) = 2 \times 470.82/96.485 = \qty{9.76}{eV}$; $D_0(\ce{N2+}) = 9.76 +
14.53 - 15.58 = \qty{8.71}{eV}$.
\textbf{19.} O orbital ocupado mais alto de \ce{N2} é ligante, abaixo do nível 2p
do átomo.
\textbf{20.} Quando o elétron retirado vem de um \omterm{def:b2:diatomic-mos:bonding}{orbital antiligante}.
\textbf{21.} A regra de Hund (número máximo de spins paralelos em orbitais
degenerados).
\textbf{22.} Os dois elétrons $\pi^*$ no mesmo orbital, com spins emparelhados.
\textbf{23.} Dois elétrons no mesmo orbital se repelem mais e perdem a
estabilização de troca dos spins paralelos.
\textbf{24.} A \omterm{def:b2:diatomic-mos:bond-order}{ordem de ligação} continua 2; todos os elétrons emparelhados:
\omterm{def:b2:diatomic-mos:magnetism}{diamagnético}.
\textbf{25.} \ce{O2+} tem \omterm{def:b2:diatomic-mos:bond-order}{ordem de ligação} $\boldsymbol{2.5}$ e energia de dissociação de
$\boldsymbol{\approx \qty{6.66}{eV}}$, contra 2 e \qty{5.12}{eV} para \ce{O2}.
\end{solution}
